\documentclass[crop=false]{standalone}
% --- ПРЕАМБУЛА ДЛЯ ПОДДЕРЖКИ РУССКОГО ЯЗЫКА И МАТЕМАТИКИ ---
\usepackage[T2A]{fontenc}
\usepackage[utf8]{inputenc}
\usepackage[russian]{babel}
\usepackage{amsmath}
\usepackage{geometry}
\usepackage{amssymb}
\usepackage{graphicx}
\usepackage{float}
\usepackage{import}
\geometry{a4paper, top=2cm, bottom=2cm, left=2cm, right=2cm}
\newcommand{\sidecomment}[1]{\marginpar{\raggedright\scriptsize\itshape #1}}
\newcommand{\iu}{\mathrm{i}}

\newcommand{\Def}{%
    \text{Def}%
    \hspace{-0.85em}% Сдвиг назад ровно на середину (подберите под шрифт)
    \raisebox{0.1ex}{/}% Черта
    \hspace{1em}% Возврат к концу слова + небольшой отступ
}

\renewcommand{\Re}{\operatorname{Re}}
\renewcommand{\Im}{\operatorname{Im}}

\ifstandalone
    \newcommand{\lecturebreak}{} % В отдельном файле лекции разрыв не нужен
\else
    \newcommand{\lecturebreak}{\clearpage} % В полной книге - начинаем с новой страницы
\fi

\newcounter{lecturenumber} % Создаем счетчик для нумерации лекций
\newcommand{\lecture}[2]{
    \lecturebreak
    \refstepcounter{lecturenumber} % Увеличиваем счетчик на 1
    \par\vspace{2em}\noindent % Отступ сверху и убираем абзацный отступ
    {\Large\bfseries % Делаем заголовок крупным и жирным
    Лекция \thelecturenumber: #1 % Выводим "Лекция N: Тема"
    \hfill % Растягиваем пробел, чтобы дата была справа
    \large\textit{#2}} % Выводим дату курсивом
    \par\vspace{1em} % Небольшой отступ снизу
    \addcontentsline{toc}{section}{Лекция \thelecturenumber: #1} % Добавляем в оглавление
    \par
}

\newcommand{\TwoCols}[2]{%
  \par\noindent % Начать с новой строки без отступа
  \begin{minipage}[t]{0.48\textwidth}
    #1
  \end{minipage}% <-- Важный '%' для удаления пробела
  \hfill % Растягивающийся пробел между колонками
  \begin{minipage}[t]{0.48\textwidth}
    #2
  \end{minipage}% <-- Важный '%' для удаления пробела
  \par%
}

\pagestyle{plain} % Включаем нумерацию страниц\

\begin{document}
\lecture{Кинематические Инварианты. Кинематические винт. Сложное движение. Углы Эйлера}{03.02}
\section{Кинематические Инварианты}
\begin{enumerate}
    \item Квадрат угловой скорости: $I_1 = (\bar \omega, \bar \omega) = \omega^2$
    \item 
    \[ \bar v_Q = \bar v_P + \bar \omega \times \overline{PQ} \implies (\bar v_Q, \bar \omega) = (\bar v_P , \bar \omega) + (\bar \omega \times \overline{PQ}, \bar \omega) \qquad (\bar \omega \times \overline{PQ}, \bar \omega) = 0\] %fixme просто подписать, что это нуль прямо в строке
    $I_2 = (\bar v_P, \bar \omega) = \text{пр}_{\bar \omega} \bar v_p \cdot w$
    \item 
    \[\bar v = \bar v_P + \bar \omega \times (\bar r - \bar r_P) = \alpha \bar \omega \qquad (\bar v_P, \bar \omega) + (\bar \omega \times (\bar r - \bar r_P), \bar \omega) = \alpha \omega^2 \quad (\bar \omega \times (\bar r - \bar r_P), \bar \omega) = 0 \quad (\bar v_P, \bar \omega) = I_2\]
    \[ \alpha = \frac{I_2}{I_1} \qquad \bar v_P + \bar \omega \times (\bar r - \bar r_P) = \frac{I_2}{I_1} \bar \omega\]
    Это уравнение называется уравнением оси винта
    \[\begin{pmatrix}
        v_{Px} \\ v_{Py} \\ v_{Pz}
    \end{pmatrix} + \begin{pmatrix}
        p \\ q \\ r
    \end{pmatrix} \times \begin{pmatrix}
        x - x_P \\ y - y_P \\ z - z_P
    \end{pmatrix} = \frac{I_2}{I_1} \begin{pmatrix}
        p \\ q \\ r
    \end{pmatrix}
    \]
    \[ 
    \begin{cases}
    v_{Px} + q (z-z_P) - r(y - y_P) = \frac{I_2}{I_1} p, \\
    v_{Py} - p (z - z_P) + r (x - x_P) = \frac{I_2}{I_1}q , \\
    v_{Pz} + p(y - y_P) - q(x - x_P) = \frac{I_2}{I_1} r
    \end{cases}
    \]
    Решение этой системы это геометрическое место точек, скорости которых коллинеарны угловой скорости % fixme что такое кинематический винт и что такое ось этого винта
\end{enumerate}
\section{Сложное движение}
\subsection{Скорость}
$r'$ дан в подвижной системе координат
\[\bar r = \bar r_{O'} + A r'\]
$\bar r'(t)$  $v$ - ? , $w$ - ?
\[\bar v = \dot{\bar r} = \frac{d}{dt}{\left(\bar r_{O'} + A \bar r \right)} = \bar v_{O'} + \dot A \bar r' + A \dot {\bar r}'  = \bar v_{O'} + \bar \omega \times (\bar r - \bar r_{O'}) + A \dot {\bar r}'\]
\[\dot A r'= \dot A A^T A r' = \dot A A^T (\bar r - \bar r_{O'}) = \Omega (\bar r - \bar r_{O'}) = \bar \omega \times (\bar r - \bar r_{O'}) \]
\[v^{\text{пер}} := \bar v_{O'} + \bar \omega \times (\bar r - \bar r_{O'}) \qquad \bar v^{\text{отн}} := A \dot {\bar r}'\]
\Def Переносная скорость – скорость с которой движется ск относительно неподвижной ск

Доказали 
\begin{theorem}[Кориолиса]
    \[\bar v = \bar v^{\text{пер}} + \bar v^{\text{отн}}\]
\end{theorem}
\subsection{Ускорение}
\[
\begin{aligned}
\bar w &= \dot {\bar v} = \dot{\bar v}_{O'} + \dot {\bar \omega} \times (\bar r - \bar r_{O'}) + \bar \omega \times \dot{\left(\bar r - \bar {r_{O'}}\right)} + \dot A \dot {\bar r}' + A \ddot{\bar r}' \\ 
&= \underbracket{\bar w_{O'} + \bar \varepsilon \times (\bar r - \bar r_{O'}) + \bar \omega \times(\bar \omega \times (\bar r - \bar r_{O'}))}_{\bar w^{\text{пер}} \quad \text{формула Ривальса}} + \bar \omega \times \underbracket{A \dot{\bar r}'}_{\bar v^{\text{отн}}} + \dot A \dot{\bar r}' + \underbracket{A \ddot{\bar r}'}_{\bar w^{\text{отн}}} \\
&\left| \dot A \dot{\bar r}' = \dot A A^T A \dot{ \bar r}' = \Omega \bar v^{\text{отн}} = \bar \omega \times v^{\text{отн}}\right| \\
&= \bar w^{\text{пер}} + \underbracket{2(\bar \omega \times v^{\text{отн}})}_{\bar w^{\text{Кор}}} + \bar w^{\text{отн}} = \bar w^{\text{пер}} + \bar w^{\text{кор}} + \bar w^{\text{отн}} 
 \end{aligned}
 \]
\begin{theorem}[Кориолиса]
    \[\bar v = \bar v^{\text{отн}} + \bar v^{\text{отн}}\]
    \[\bar w = \bar w^{\text{пер}} + \bar w^{\text{Кор}} + \bar w^{\text{отн}} \]
    \[\bar w^{\text{Кор}} = 2 (\bar w \times \bar v^{\text{отн}})\]
\end{theorem}
\section{Сложное движение}
\[A = A_2 A_1 = A_1 A_2'\]
\[\bar r = A \bar r_0 = A_1 A_2' \bar r_0\]
\[\bar v = \Omega \bar r = \bar \omega \times \bar r\]
\begin{proof}\[
\begin{aligned}
\bar v &= \dot A_1  A_2 \bar r_0 + A_1 \dot A_2' \bar r_0 = \underbracket{\dot A_1 A_1^T}_{\Omega_1} \underbracket{A_1 A_2' \bar r_0}_r + A_1 \underbracket{\dot A_2' A_2'^T}_{\Omega_2'} A_2' r_0 = \\
&= \Omega_1 \bar r + A_1 \Omega_2' A_2' \bar r_0 = \Omega_1 \bar r + \Omega_2 \bar r = \bar \omega \times \bar r + \bar \omega_2 \times \bar r % fixme если непонятно, есть фото в телефоне
\end{aligned}
\]
\end{proof}
\[\bar \varepsilon = \frac{d}{dt} \bar \omega = \dot{\bar \omega} = \frac{d}{dt}{\left(\bar \omega_1 + \bar \omega_2\right)} = \frac{d}{dt} (\bar \omega_1 + A_1 \bar \omega_2') = \bar \varepsilon_1' + \dot A_1 \bar \omega_2' + A_1 \bar \omega_2'  \qquad A_1 \dot{\bar {\omega}}_2' = A_1 \bar \varepsilon_2' = \bar \varepsilon_2  \]
\[\bar \varepsilon = \bar \varepsilon_1 + \bar \omega_1 \times \bar \omega_2 + \bar \varepsilon_2 \qquad \dot A_1 \bar \omega_2' = \dot A_1 A_1^T A_1 \bar \omega_2' = \Omega \bar \omega_2' = \bar \omega_1 \times \bar \omega_2\]
\[\boxed{\bar \varepsilon = \bar \varepsilon_1 + \bar \varepsilon_2 + \bar \omega_1 \times \bar \omega_2}\]
\end{document}